% Generated from validated Article Draft Result. Do not hand-edit; revise the structured draft instead.
\section{Multimodal \& Generation — 生成・理解・groundingが別々に伸びる}
\label{pkg:multimodal-generation-march}
\noindent\textbf{LongCat-Next、Unify-Agent、DreamLite、Gen-Searcher。3月のmultimodal領域は、共通token空間、world-grounded image synthesis、on-device generation、search-augmented generationという異なる方向へ同時に広がった。}\autocite{src-197775677c0d,src-abece2adf5a7,src-ac36472a8f09,src-b8369fd2c988}

\subsection{MultimodalをAgentの一部だけで読まない}

3月のmultimodal/generation研究には、modalitiesを共通表現へ寄せるmodeling、端末上で生成と編集を両立するdeployment、外部検索で生成をgroundするpipelineが並ぶ。Agent的な構成もあるが、すべてをagent narrativeへ吸収すると、generation/understandingそのものの設計変化を見失う。\autocite{src-197775677c0d,src-abece2adf5a7,src-ac36472a8f09,src-b8369fd2c988}


\subsection{LongCat-Next: modalityをshared discrete spaceへ}

LongCat-NextはauthorsがDiNA（Discrete Native Autoregressive）を提案し、multimodal informationをshared discrete spaceで扱ってmodalities横断のautoregressive modelingを統一する方向を示した。ここで重要なのは、画像生成だけの話ではなく、generationとunderstandingを同一の表現・生成原理へ寄せようとする設計である。\autocite{src-ac36472a8f09}


\subsection{Image generationへsearchとgroundingを持ち込む}

Unify-Agentはauthorsがprompt understanding、multimodal evidence search、grounded recaptioning、final synthesisから成るworld-grounded image synthesis pipelineとして提案した。Gen-Searcherもsearch-augmented image generation agentとして、multi-hop reasoningとsearchでtext knowledgeやreference imageを収集する構成を示す。両者はAgent技術との接点を持つが、目的は画像生成を外部evidenceへgroundすることにあり、一般的なagent runtimeとは区別して読むべきだ。\autocite{src-197775677c0d,src-abece2adf5a7}


\subsection{DreamLite: on-deviceでも生成と編集を統合する}

DreamLiteはauthorsが0.39Bのcompact unified on-device diffusion modelとして、text-to-image generationとtext-guided image editingを単一networkで扱う構成を提案した。frontier-scale modelとは別の方向から、generation capabilityをdevice側へ近づける研究である。\autocite{src-b8369fd2c988}


3月のmultimodal領域は、unification、grounding、deploymentという複数の設計軸を持っていた。後から見ればagentic synthesisへの接続も目立つが、当月のEvidenceだけから見ても独立した主要テーマとして扱うだけの厚みがある。\autocite{src-197775677c0d,src-abece2adf5a7,src-ac36472a8f09,src-b8369fd2c988}


\begin{claimboundary}
Claim Boundary: 各architectureと評価結果はauthorsの報告であり、本誌はbenchmarkを独立再現していない。March chronologyは初回投稿を基準とし、Unify-Agentのv2、Gen-Searcherのv3など後日revisionの内容を3月時点へ遡及させない。LongCat-Nextについても3月初回投稿として扱い、後から得た知見で当月の記述を膨らませない。\autocite{src-197775677c0d,src-abece2adf5a7,src-ac36472a8f09,src-b8369fd2c988}
\end{claimboundary}
